% TOP_PROBLEM: {"id": 30006701, "title": "Finiteness of the Tate-Shafarevich Group for Elliptic Curves", "category_id": 1, "category": {"id": 1, "name": "number_theory", "display_name": "Number Theory", "description": "Properties of integers, prime numbers, Diophantine equations.", "slug": "number-theory", "order_index": 1, "created_at": "2026-07-31T15:26:25.670Z"}, "status": "open"}
% FIRST_POSTED: 2026-09-27
% !TeX program = lualatex
% Compile twice: lualatex 30006701.tex
% Original entry retained; research subsection and [E] references added.
% No external bibliography, image, data, or style file is required.
\documentclass[11pt,a4paper]{article}
\usepackage[margin=24mm,headheight=14pt]{geometry}
\usepackage{fontspec}
\setmainfont{Latin Modern Roman}
\setsansfont{Latin Modern Sans}
\setmonofont{Latin Modern Mono}
\usepackage{amsmath,amssymb,mathtools,amsthm}
\usepackage{unicode-math}
\setmathfont{Latin Modern Math}
\usepackage{microtype}
\usepackage{enumitem,xurl,needspace}
\usepackage[dvipsnames]{xcolor}
\usepackage{hyperref}
\hypersetup{unicode=true,colorlinks=true,linkcolor=MidnightBlue,
 urlcolor=MidnightBlue,bookmarksnumbered=true,
 pdftitle={Research notebook - Problem 30006701},
 pdfauthor={Research attempt; not independently refereed}}
\usepackage{bookmark}
\usepackage{titlesec}
\titleformat{\section}{\Large\bfseries\raggedright}{\thesection}{0.7em}{}
\usepackage{fancyhdr}
\pagestyle{fancy}\fancyhf{}
\fancyhead[L]{\small\sffamily Research notebook}
\fancyhead[R]{\small\sffamily Problem 30006701}
\fancyfoot[L]{\scriptsize 27 September 2026}
\fancyfoot[C]{\thepage}
\fancyfoot[R]{\scriptsize Unrefereed research attempt}
\setlength{\parindent}{0pt}
\setlength{\parskip}{5pt plus 1pt}
\setlength{\emergencystretch}{3em}
\clubpenalty=10000
\widowpenalty=10000
\displaywidowpenalty=10000
\setcounter{secnumdepth}{1}
\setlist[itemize]{leftmargin=3em,itemsep=5pt,topsep=4pt,parsep=0pt}
\allowdisplaybreaks
\newcommand{\N}{\mathbb{N}}
\newcommand{\Z}{\mathbb{Z}}
\newcommand{\Q}{\mathbb{Q}}
\newcommand{\R}{\mathbb{R}}
\newcommand{\C}{\mathbb{C}}
\newcommand{\F}{\mathbb{F}}
\newcommand{\A}{\mathbb{A}}
\newcommand{\PP}{\mathbb P}
\newcommand{\E}{\mathbb{E}}
\newcommand{\eps}{\varepsilon}
\newcommand{\dd}{\,\mathrm d}
\newcommand{\sref}[2]{\hyperlink{source:#1:#2}{[S#2]}}
\newcommand{\eref}[2]{\hyperlink{extra:#1:#2}{[E#2]}}
\newif\ifshowcatalogue
\showcataloguetrue
\newcommand{\cataloguescope}[1]{\ifshowcatalogue
 \paragraph{Exact scope in the supplied catalogue.}
 \begin{quote}\small #1\end{quote}\fi}
\newtheorem{notebooklemma}{Lemma}[section]
\newtheorem{notebookproposition}[notebooklemma]{Proposition}
\numberwithin{equation}{section}
\begin{document}
\setcounter{section}{0}
% ================================================================
% problemId: 30006701
% Canonical problem ID: 30006701
% exactTarget SHA-256: 2653bdb6ffc8f0769f408a872e5727fb611765a9bc0b7c3e3e0376beeec87885
\clearpage
\section*{\texorpdfstring{Finiteness of the Tate-Shafarevich Group for Elliptic Curves}{Finiteness of the Tate-Shafarevich Group for Elliptic Curves}}\label{30006701}
\subsection{Definitions and mathematical statement}
For an elliptic curve $E/K$, define the Tate--Shafarevich group by
\[
 \operatorname{Sha}(E/K)=\ker\!\left(H^1(K,E)\longrightarrow\prod_v H^1(K_v,E)\right),
\]
where $v$ runs over all places of the number field $K$ and $H^1$ is Galois cohomology. Its elements classify torsors under $E$ having a point over every completion $K_v$. The conjecture is $\#\operatorname{Sha}(E/K)<\infty$ for every $E/K$. It concerns the whole group, not only each fixed finite-torsion subgroup or a single primary component.
\par\smallskip\textit{Formulation and concept references:} \sref{30006701}{1}.
\cataloguescope{Prove that the Tate-Shafarevich group Sha(E/K) is finite for every elliptic curve E defined over a number field K.}
\subsection{Short English statement}
Are there only finitely many locally soluble torsor classes for every elliptic curve over every number field?
\subsection{Research attempt}

\paragraph{Current frontier.}
For each fixed integer $m$, descent gives a finite $m$-Selmer group and hence finite $\operatorname{Sha}(E/K)[m]$. That does not imply finiteness of the whole group. The classical Gross--Zagier--Kolyvagin conclusion applies, in particular, to elliptic curves over $\Q$ of \emph{analytic} rank at most one; algebraic rank at most one alone is not the hypothesis being used \eref{30006701}{1}. The Cassels--Tate pairing for an elliptic curve is alternating, with kernel the maximal divisible subgroup; this also does not on its own remove that subgroup \eref{30006701}{2}. Banwait's September 2026 manuscript supplies primewise criteria and finite computations for rank-two CM curves, and explicitly does not infer whole-group finiteness from them \eref{30006701}{3}. We do not assume its new claims in the proof below.

\paragraph{Attack route.}
An Euler-system argument could eliminate a fixed $p$-primary divisible part, but a separate assertion excluding infinitely many nonzero primary components would still be needed. A direct descent computation might bound a Selmer group at a chosen level, but its liftable classes must be distinguished from its genuinely new obstructions. The route explored here synchronizes \emph{all} primes and arbitrarily deep powers of each prime by using factorial descent levels. It produces a quantitative criterion with an explicit threshold, and an abstract sharpness example shows how much can be concluded from the standard group-theoretic properties alone.

\paragraph{Attempt: separate the two ways the group could be infinite.}
Fix the curve $E/K$ once and for all and write $A=\operatorname{Sha}(E/K)$. Standard Galois cohomology and the Kummer sequence give
\begin{equation}\label{r59:kummer}
 0\longrightarrow E(K)/mE(K)\longrightarrow
       \operatorname{Sel}_m(E/K)\longrightarrow A[m]\longrightarrow0.
\end{equation}
Here $A$ is torsion and $A[m]$ is finite for each $m$ \eref{30006701}{1}, Chapter IV, Sections 1--3. The following deduction explains why these finite sets do not automatically exhaust a finite group.

\begin{notebooklemma}[An infinite primary part contains a divisible chain]\label{r59:prufer}
If a torsion abelian group $A$ has finite $A[p^a]$ for every $a\geq1$, then either $A[p^\infty]$ is finite or it contains a subgroup isomorphic to $\Q_p/\Z_p$.
\end{notebooklemma}
\begin{proof}
If the orders in $A[p^\infty]$ were bounded by $p^a$, the whole primary group would be contained in the finite group $A[p^a]$. Otherwise there are elements of arbitrarily large exact $p$-power order. Form a rooted tree of finite strings $(x_1,\ldots,x_a)$ with $x_i$ of exact order $p^i$ and $px_{i+1}=x_i$. Each level is finite, because $A[p^a]$ is finite, and there are strings of every length by taking multiples of an element of order $p^a$.

A finitely branching tree with arbitrarily long branches has an infinite branch: choose successively a child with descendants at arbitrarily large levels. The union of the nested cyclic groups along that branch is the Pr\"ufer $p$-group, isomorphic to $\Q_p/\Z_p$. In particular it contributes at least $p^a$ elements to $A[p^a]$ for every $a$.
\end{proof}

Next use the pairing, but only where it is justified. If $A[p^\infty]$ is finite, it has zero intersection with the maximal divisible subgroup. Elements of distinct primary parts pair to zero in $\Q/\Z$, since the value is killed by two coprime powers. Nondegeneracy modulo the divisible subgroup therefore makes the restriction to this finite primary part nondegenerate. Because the canonical principal polarization of an elliptic curve comes from its rational origin, the pairing here is alternating, including at $p=2$ \eref{30006701}{2}.

\begin{notebooklemma}[The first layer costs at least two dimensions]\label{r59:two}
If $A[p^\infty]$ is finite and nonzero, then $|A[p]|\geq p^2$.
\end{notebooklemma}
\begin{proof}
Let $G=A[p^\infty]$ and choose $x\in G$ of maximum order $p^a$. The map $G\to\operatorname{Hom}(G,\Q/\Z)$ from its nondegenerate pairing is injective and hence bijective, since both finite groups have the same size. The character paired with $x$ has exact order $p^a$, so some $y$ pairs with $x$ to an element of exact order $p^a$. Necessarily $y$ also has order $p^a$.

If $ux+vy=0$, pairing successively with $x$ and $y$ shows that $u$ and $v$ are divisible by $p^a$, because the pairing is alternating. Thus $x,y$ generate a subgroup isomorphic to $(\Z/p^a\Z)^2$, whose $p$-torsion already has $p^2$ elements.
\end{proof}
The distinction between the two lemmas is essential: Lemma~\ref{r59:two} is not being applied to a primary group that might still have a divisible part.

\paragraph{Attempt: a synchronized Selmer-growth criterion.}
For an integer $N\geq2$, define the positive integer
\begin{equation}\label{r59:defect}
 D_E(N)=\frac{|\operatorname{Sel}_{N!}(E/K)|}
                   {|E(K)/N!E(K)|}=|A[N!]|.
\end{equation}
This exact equality, rather than an asymptotic formula, follows from \eqref{r59:kummer}. If $E(K)\simeq\Z^r\oplus T$ with $T$ finite, its denominator is
\[
 |E(K)/N!E(K)|=(N!)^r\,|T/(N!)T|,
\]
which equals $(N!)^r|T|$ once $N!$ is divisible by the exponent of $T$. No unknown Mordell--Weil torsion factor is discarded.

\Needspace{10\baselineskip}
\begin{notebookproposition}[A strict quadratic threshold]\label{r59:threshold}
The following are equivalent for the fixed elliptic curve $E/K$:
\[
 \#A<\infty;
 \qquad
 \limsup_{k\to\infty}
   \frac{\log D_E(2^k)}{\log(2^k)}<2.
\]
Equivalently, it suffices to prove that there is an $\eta>0$, allowed to depend on $E/K$, such that
\begin{equation}\label{r59:missing}
 D_E(2^k)\leq(2^k)^{2-\eta}
                   \quad\text{for every sufficiently large }k.
\end{equation}
\end{notebookproposition}
\begin{proof}
If $A$ is finite, then $A[N!]=A$ for all sufficiently large $N$, so the limit is zero.

Conversely suppose \eqref{r59:missing}. If $A[p^\infty]$ were infinite for some fixed prime $p$, Lemma~\ref{r59:prufer} would give
\[
 D_E(N)\geq p^{v_p(N!)}\geq p^{\lfloor N/p\rfloor}.
\]
Along dyadic $N$ this eventually exceeds every fixed power of $N$, a contradiction. Thus every primary part is finite.

If infinitely many of them were nonzero, choose such a prime $p$ arbitrarily large and let $N$ be the least power of two with $N\geq p$. Then $p\leq N<2p$ and $p\mid N!$. By Lemma~\ref{r59:two},
\[
 D_E(N)\geq|A[p]|\geq p^2>\frac{N^2}{4}.
\]
For $N^\eta>4$ this contradicts \eqref{r59:missing}. Hence only finitely many primary parts are nonzero, and all are finite. A torsion abelian group is the direct sum of its primary parts, so $A$ is finite. Finally a limsup strictly less than two is equivalent to an eventual bound with a fixed positive gap below two, by the definition of limsup.
\end{proof}

This criterion handles both the depth of a fixed primary component and the occurrence of entirely new primes. In particular it is not the incorrect inference that finiteness of all individual $A[p^a]$, or even of all $A[p^\infty]$, settles the target. It also shows exactly why an exponent threshold of two appears when the alternating pairing is available.

\paragraph{Attempt: actively testing the threshold by an infinite model.}
Choose increasing odd primes $p_j$ so rapidly that
\begin{equation}\label{r59:lacunary}
 \sum_{i<j}\log p_i=o(\log p_j).
\end{equation}
Such a sequence can be chosen inductively using only the existence of arbitrarily large primes. Let
\[
 A_0=\bigoplus_{j\geq1}(\Z/p_j\Z)^2,
\]
and equip the $j$th summand with
\[
 \langle(a,b),(c,d)\rangle_j=\frac{ad-bc}{p_j}\pmod\Z.
\]
Sum the pairings across components. Each group element has finite support, so this is a well-defined alternating pairing with zero left and right kernels. Every fixed $A_0[m]$ is finite and every primary part is finite, but $A_0$ itself is infinite and has no nonzero divisible subgroup.

For its factorial defect,
\begin{equation}\label{r59:toydefect}
 D_{A_0}(N):=|A_0[N!]|=\prod_{p_j\leq N}p_j^2.
\end{equation}
When $p_j\leq N<p_{j+1}$,
\[
 \frac{\log D_{A_0}(N)}{\log N}
 \leq\frac{2\sum_{i\leq j}\log p_i}{\log p_j}=2+o(1).
\]
At the least dyadic integer $N_j\geq p_j$, one has $N_j<2p_j<p_{j+1}$ for large $j$, and the same quotient tends to two by \eqref{r59:lacunary}. Therefore
\begin{equation}\label{r59:sharp}
 \limsup_{k\to\infty}\frac{\log D_{A_0}(2^k)}{\log(2^k)}=2,
 \qquad \#A_0=\infty.
\end{equation}
This proves that the strict threshold in Proposition~\ref{r59:threshold} cannot be replaced by a non-strict one using only these abstract torsion and pairing properties. It does \emph{not} construct a curve having $A_0$ as its Tate--Shafarevich group. Such an arithmetic realization would be an actual counterexample and is wholly missing.

\paragraph{Stress test: why a pairing computation does not close the gap.}
Even a perfect pairing on a finite group may become degenerate on its first torsion layer. On $G=(\Z/p^2\Z)^2$, use the perfect alternating pairing
\[
 \langle(a,b),(c,d)\rangle=\frac{ad-bc}{p^2}\pmod\Z.
\]
The restriction to $G[p]=pG$ is identically zero, since its numerator is divisible by $p^2$. Therefore zero pairings among $p$-torsion classes do not certify that those classes come from Mordell--Weil points, and do not identify the radical needed at deeper descent levels. The more general Selmer-pairing formalism explicitly tracks images from other Selmer groups \eref{30006701}{4}. Likewise $\Q_p/\Z_p$ has finite first layer but an infinite primary tower, so bounding only that first layer misses a different obstruction.

The positive direction now requires an arithmetic estimate of the form \eqref{r59:missing}, uniformly over sufficiently large factorial levels for each fixed curve. Nothing above proves such an estimate. Its exact equivalent expression is
\[
 |\operatorname{Sel}_{N!}(E/K)|
       \leq |E(K)/N!E(K)|\,N^{2-\eta}
                 \quad(N=2^k\text{ sufficiently large}).
\]
Standard finiteness of each descent group gives no such uniform growth control. Replacing its right side by an uncontrolled level-dependent constant would remove all content. Neither a finite list of primes satisfying a CM criterion nor nonvanishing of a regulator at selected primes supplies this estimate \eref{30006701}{3}. The full number field $K$, the fixed curve, the arbitrarily deep powers, and the unbounded prime range are all retained.

The accompanying exact script checks that the displayed pairing on $(\Z/p^2\Z)^2$ vanishes on its $p$-torsion for $p=2,3,5,7$. It also checks the factorial valuation bound and the dyadic inequalities used in the proof; for $N=8,16,32,64$, the values of $v_3(N!)$ are $2,6,14,30$. These computations do not evaluate $D_E(N)$ for an elliptic curve. In particular no numerical evidence for \eqref{r59:missing} is claimed.

\paragraph{Outcome.}
\textbf{Outcome: Conditional reduction.}

The argument proves a synchronized growth criterion for whole-group finiteness, and an infinite abstract pairing group proves sharpness of its strict exponent threshold at the group-theoretic level. It supplies neither the required arithmetic Selmer estimate for a new curve nor a realization of the abstract infinite group as a Tate--Shafarevich group. The reduction is therefore not a finiteness theorem for any new elliptic-curve case. The all-curve conclusion would follow from \eqref{r59:missing} for every $E/K$, with constants allowed to depend on that curve and field. Novelty and correctness of these intermediate deductions have not received independent expert verification.

\subsection{Sources}
\begin{itemize}\raggedright
\item[\textnormal{[S1]}]\hypertarget{source:30006701:1}{} J. H. Silverman, The Arithmetic of Elliptic Curves, GTM 106, Springer, 2009.
\url{https://doi.org/10.1007/978-0-387-09494-6}.
{\small\textit{Catalogue formulation source.}}
\item[\textnormal{[S2]}]\hypertarget{source:30006701:2}{} Rank stability in quadratic extensions and Hilbert’s tenth problem for the ring of integers of a number field.
\url{https://link.springer.com/article/10.1007/s00222-025-01392-3}.
{\small\textit{Catalogue research/status reference; not a proof endorsement.}}
\item[\textnormal{[S3]}]\hypertarget{source:30006701:3}{} Finiteness of the Tate-Shafarevich group over function fields for groups of multiplicative type.
\url{https://arxiv.org/abs/2607.15372}.
{\small\textit{Catalogue research/status reference; not a proof endorsement.}}
\item[\textnormal{[S4]}]\hypertarget{source:30006701:4}{} Second derivatives of p-adic L-functions and the Shafarevich–Tate group of rank-two CM elliptic curves.
\url{https://arxiv.org/abs/2609.08431}.
{\small\textit{Catalogue research/status reference; not a proof endorsement.}}

% Research references added for this notebook.
\item[\textnormal{[E1]}]\hypertarget{extra:30006701:1}{} J. S. Milne, \emph{Elliptic Curves}, second edition, World Scientific (2021), Chapter IV, Sections 1--5 and 10, especially equation (29), Theorem 3.1, and the analytic-rank-at-most-one conclusion in Section 10.
\url{https://www.jmilne.org/math/Books/EC2.pdf}.
{\small\textit{Kummer exact sequence, finite Selmer groups, torsion and divisible-primary issues, and the correctly qualified classical finiteness theorem. Consulted 27 September 2026.}}
\item[\textnormal{[E2]}]\hypertarget{extra:30006701:2}{} Bjorn Poonen and Michael Stoll, \emph{The Cassels--Tate pairing on polarized abelian varieties}, Annals of Mathematics (2) 150 (1999), no. 3, 1109--1149; arXiv:math/9911267, Introduction and the discussion of the elliptic-curve alternating case.
\url{https://arxiv.org/pdf/math/9911267}.
{\small\textit{Nondegeneracy modulo the maximal divisible subgroup; alternating pairing for elliptic curves, including its relevance at 2. Consulted 27 September 2026.}}
\item[\textnormal{[E3]}]\hypertarget{extra:30006701:3}{} Barinder S. Banwait, \emph{Second derivatives of $p$-adic $L$-functions and the Shafarevich--Tate group of rank-two CM elliptic curves}, arXiv:2609.08431v1 (8 September 2026), Section 1, Theorems A--B, and the stated limitations of the computations and formalisation.
\url{https://arxiv.org/html/2609.08431v1}.
{\small\textit{Recent primewise criteria and finite computational range; no whole-group finiteness is inferred. The notebook does not use its new claims as a theorem. Consulted 27 September 2026.}}
\item[\textnormal{[E4]}]\hypertarget{extra:30006701:4}{} Adam Morgan and Alexander Smith, \emph{The Cassels--Tate pairing for finite Galois modules}, arXiv:2103.08530v3 (8 February 2023), abstract and the main pairing construction and kernel statement.
\url{https://arxiv.org/html/2103.08530v3}.
{\small\textit{Selmer pairing kernels are images from appropriate other Selmer groups, not automatically just Mordell--Weil images. Consulted 27 September 2026.}}
\end{itemize}

\end{document}
